\documentclass[12pt]{article}
\usepackage[utf8]{inputenc}
\usepackage[T1]{fontenc}
\usepackage{hyperref}
\title{Komo Language Resources for K-12 Curriculum Development}
\author{Compiled by Claude Sonnet 5.5 (claude-sonnet-5-5) for Andr\'as Kornai}
\date{}
\begin{document}
\maketitle

\section{Introduction}
Komo (ISO 639-3 xom; native names Tta Komo and Go-kwom) belongs to the Koman group. It is spoken in Ethiopia (Benishangul-Gumuz and Gambela) and Sudan. Reported figures are 8,500 speakers in Ethiopia (2007 census) and 10,000 in Sudan (1979). Sources describe it as greatly understudied. The search found no dictionary. Be careful not to confuse it with the related Gwama (Kwama) language.

\section{Primary Educational Resources}
None were found.

\section{Linguistic and Reference Materials}
\begin{itemize}
\item Wikipedia, Komo language: \url{https://en.wikipedia.org/wiki/Komo_language}
\item Wikipedia, Koman languages: \url{https://en.wikipedia.org/wiki/Koman_languages}
\item Gwama introduction (SOAS), for the related language: \url{https://glocal.soas.ac.uk/afala2023-91/}
\item Gwama Webonary introduction, for the related language, not Komo: \url{https://www.webonary.work/gwama/overview/introduction/}
\end{itemize}

\section{Religious and Community Resources}
None were found.

\section{Digital and Audio Resources}
\begin{itemize}
\item Omniglot page with an alphabet chart and sample text: \url{https://www.omniglot.com/writing/komo.htm}
\end{itemize}

\section{Curriculum Development Context}
Few Komo-specific resources appear online. Omniglot shows a writing system in use, but no school materials were found. The Gwama resources are for a different language and are listed only as regional context.
\end{document}